\section{The exact auxiliary-twist obstruction at the fixed prime \(269\)}
\label{sec:fixed-cubic-twist-admissibility}

The source permits more characters in its theta-data inclusion
than in its quantitative good-twist theorem.  The difference is
decisive for a fixed cubic character, even after a suitable CM
host and principal polarization have been constructed.

\begin{proposition}[AC-G40: fixed cubic inertia fails the literal good-twist conditions]
\label{AC-G40}
Let $F$ be a number field containing $K=\mathbb Q(\omega)$
and unramified at $p=269$.  Put
$F^{\mathrm c}=F(\alpha)$, $\alpha^3=p$, and let
$\chi:G_F\to\mu_3$ be either faithful character of this
cyclic extension.  Write $\mathcal P_p$ for all primes
of $F$ above $p$.  Let $\eta$ be a finite character
unramified at every member of $\mathcal P_p$; it may
be trivial.  Put $\phi=\chi\eta$.

For every $P\in\mathcal P_p$, the inertia image of
$\phi$ is exactly $\mu_3$, and its absolute residue
degree $f(P/p)$ is even.  Thus any set $Q$ satisfying
the source's \texttt{RamExactly} condition for $\phi$
contains every member of $\mathcal P_p$.
In particular $\phi$ cannot satisfy the literal
\texttt{BlochKatoReduction.Gen.PsiOk} or
\texttt{IntegralHeckeForms.GoodTw}, because
\texttt{PsiOk} requires absolute residue degree one
at every prime in $Q$.

There are two additional distinctions by coefficient prime.
If $\ell\ne3$, the source requirement that all values
of $\phi$ lie in the cyclic group generated by a
primitive $\ell^{k'}$-th root already fails on inertia
at $p$.  This includes $\ell=p=269$; at that coefficient
prime the requirement \texttt{TrivAbove} also fails.
At $\ell=3$, the precise source-normalized tame value is
\[
                  \operatorname{tameVal}_3(\phi,P)=\frac12
                       \qquad(P\in\mathcal P_p).              \tag{G40.1}
\]
Consequently its \texttt{partD} expression satisfies
\[
       t\sum_{P\in Q}\operatorname{tameVal}_3(\phi,P)
                         \geq\frac{t\,\#\mathcal P_p}{2}.
                                                               \tag{G40.2}
\]
It cannot be at most the literal choice
$\texttt{twEps}=3^{-2}=1/9$ when $t\geq1$.
For $\eta=1$, $Q=\mathcal P_p$ is the exact conductor
support, and equality holds in (G40.2).
\end{proposition}
\begin{proof}
\proofstep{1. Residue degree and every ramified prime.}
The polynomial $X^2+X+1$ is irreducible modulo $269$.
A root would be a nontrivial cube root of unity in
$\mathbb F_{269}^{\times}$, whose order is $268$,
not divisible by $3$; $1$ is not a root since
$3\ne0$ in that field.  Therefore $269$ is inert
in $K$.  For any $P$ above it, the inclusion of
integer rings induces
$\mathbb F_{269^2}\hookrightarrow\mathcal O_F/P$.
The residue-degree tower formula makes $f(P/p)$
divisible by two.

Unramifiedness of $F$ at $p$ gives $v_P(p)=1$.
Hence $X^3-p$ is Eisenstein over $F_P$.  It defines
a totally ramified degree-three extension.  Since
$F_P$ contains $\mu_3$ and $p\ne3$, this extension
is cyclic and tame, with its full group equal to
the inertia image.  In particular $F^{\mathrm c}/F$
has degree three.  The inertia restriction of $\eta$
is trivial, so $\phi|_{I_P}=\chi|_{I_P}$ has image
$\mu_3$ at every $P\in\mathcal P_p$.

The source definition is an equivalence, not an
upper bound on ramification:
\[
 P\in Q\quad\Longleftrightarrow\quad
        \exists\sigma\in I_P,\ \phi(\sigma)\ne1.
                                                               \tag{G40.3}
\]
This is \texttt{RamExactly} in \texttt{DataN.lean}.
Thus every such $P$ belongs to $Q$, and its residue
degree contradicts the degree-one clause of
\texttt{PsiOk}.  Its definition of \texttt{GoodTw}
includes \texttt{PsiOk} as its first conjunct.

\proofstep{2. The coefficient-prime distinction.}
If all values of $\phi$ were powers of a primitive
$\ell^{k'}$-th root, every inertia image would have
$\ell$-power order.  Its order-three image at $p$
therefore forces $\ell=3$.  At $\ell=p$, the source
also asks that $\phi$ be trivial on the entire
decomposition group at every prime above $\ell$;
nontrivial inertia directly contradicts that.
There is no excluded coefficient-prime exception
which repairs the proposed fixed-character insertion.

\proofstep{3. The exact tame cost at three.}
Let $\mathcal R$ be any coefficient DVR of residue
characteristic three containing the character values,
with normalized integral valuation $v_{\mathcal R}$.
For a nontrivial cubic root $\zeta$,
\[
       (\zeta-1)^2=-3\zeta,\qquad
       \frac{v_{\mathcal R}(\zeta-1)}
            {v_{\mathcal R}(3)}=\frac12.
\]
The source \texttt{nval} is exactly this quotient.
Its \texttt{tameVal} is the infimum of these values
over inertia elements on which the character is
nontrivial.  At $P\in\mathcal P_p$ that set of values
is the nonempty singleton $\{1/2\}$, proving (G40.1).
At any other $P\in Q$, (G40.3) makes the defining
set nonempty and all its values are nonnegative.
This proves (G40.2).  Since $\#\mathcal P_p\geq1$,
$t\geq1$ gives a lower bound $1/2>1/9$.

\proofstep{4. Exact support and local triviality at three.}
For $\chi$ alone there is no ramification away from
$3p$: where $3p$ is a unit, the algebra defined by
$X^3-p$ is etale, since its root and derivative are
units.  At three, $-269=1-270\in1+27\mathbb Z_3$.
The convergent $3$-adic logarithm and exponential
give
\[
       \beta=\exp\!\left(\frac13\log(-269)\right)
                   \in1+9\mathbb Z_3,\qquad
       (-\beta)^3=269.
\]
Thus $X^3-p$ splits in every $F_v$ above three,
which already contains $\mu_3$; the entire local
decomposition action of $\chi$ is trivial there.
The exact support is therefore $\mathcal P_p$.
The previous inertia computation gives equality
in (G40.2) for this character.
\end{proof}

\begin{corollary}[AC-G40-inclusion: the broader theta condition remains distinct]
\label{AC-G40-inclusion}
Suppose in addition that $F$ is a CM host with every
prime in $\mathcal P_p$ split over its maximal real
subfield and disjoint from the source set $S_0$.
At coefficient three, the fixed character $\chi$
can satisfy the source condition \texttt{TwOK},
after choosing its cubic-root coefficient extension.
Consequently a supplied instance of \texttt{admT4}
does give the corresponding theta-lattice inclusion.
The available good-twist theorem therefore does
not supply its quantitative \texttt{GoodTw} conclusion
for this prescribed character.  No \texttt{IdxAt}
bound for it is established by the inclusion.
\end{corollary}
\begin{proof}
The character has finite image, hence open kernel,
and takes values in the powers of a primitive cube
root.  Step 4 proves \texttt{TrivAbove} at three,
and the proposition proves exact ramification set
$\mathcal P_p$.  Complex conjugation fixes the
real root $\alpha$ and inverts $\omega$, so
$\chi(cgc^{-1})=\chi(g)^{-1}$; this is the source's
\texttt{AntiHom} condition.  The assumed splitness
and disjointness supply the remaining clauses of
\texttt{TwOK}.  Its definition contains no
absolute residue-degree-one restriction.

If the character is required in the ray-character
notation, \texttt{NTw.art\_univ} applied to its finite
group $\mu_3$ factors it through the source
\texttt{RCG} for $Q=\mathcal P_p$, since $Q$ is
away from three.  The definition of \texttt{admT4}
then yields the inclusion
\[
       \texttt{ThPsi}(\chi)
           \ \subset\
       \texttt{thP}\bigl(\texttt{MtLat}(\chi)\bigr)
                                                               \tag{G40.4}
\]
for the admitted theta data.
The preceding proposition disproves the additional
\texttt{PsiOk} and small-\texttt{partD} conditions
for this fixed character, so the broader inclusion
does not instantiate that good-twist theorem.
\end{proof}

\begin{proposition}[AC-G40-avoidance: fresh auxiliary choices can avoid \(269\)]
\label{AC-G40-avoidance}
At a source \texttt{TwistTheoremAtIndex} instance,
whenever its preceding hypotheses hold, the
existential auxiliary character can be selected
unramified at all primes above $269$ while retaining
the same prescribed finite avoidance set and the
same quantitative allowance.  This assertion
does not prescribe that auxiliary character to
be $\chi_{269}$ or $\chi_{269}$ times a chosen character.
\end{proposition}
\begin{proof}
The source theorem quantifies over every finite
avoidance set $A$ before asserting the existence
of $Q,k',\psi'$.  Apply it to
$A^+=A\cup\mathcal P_p$.  The first conjunct of
\texttt{GoodTw} is \texttt{PsiOk} with avoidance set
$\texttt{levs}\cup A^+$.  Its disjointness clause
therefore gives $Q\cap\mathcal P_p=\varnothing$;
\texttt{RamExactly} makes the selected character
unramified there.  Every other conjunct of
\texttt{GoodTw} is independent of $A$.
The same disjointness implies the weaker clause
with $\texttt{levs}\cup A$, so this very witness
is good for the originally prescribed set and
unchanged allowance.  This is an assertion at
the fresh existential-selection stage, not a
modification of an already fixed \texttt{NChoice}.
\end{proof}

\paragraph{Quantitative scope.}
For an unramified number field containing $K$,
every residue degree above $269$ is at least two,
so $\#\mathcal P_p\leq[F:\mathbb Q]/2$.
The exact new local term for $\chi$ is at most
$t[F:\mathbb Q]/4$.  On the audited growing hosts,
$t\leq x^{401}$ and
$[F:\mathbb Q]=O_\varepsilon(x^{2+8\varepsilon})$,
so this term is polynomial and is
$o(\exp(x^\varepsilon/2))$.  This only measures the
displayed \texttt{partD} expression; it does not
extend a theorem whose residue-degree hypothesis
fails, nor supply a prescribed-character theta
index estimate.  A theorem allowing fixed inertial
support with an explicit extra cost would be a
strictly broader input than the literal source one.

\paragraph{Exact inspected source locators.}
All paths below are in the pinned source revision
listed in the evidence manifest.
In
\path{Formalism/SharpRestriction/LevelLetterCocycles/DataN.lean},
the relevant declarations are \texttt{RamExactly},
\texttt{NTw.art\_univ}, and \texttt{NTw.psiL\_l}.
In \path{BlochKatoPhases.lean}, they are
\texttt{Gen.nval}, \texttt{Gen.tameVal},
\texttt{Gen.TrivAbove}, \texttt{Gen.PsiOk}
(definition beginning at line 289), and
\texttt{Gen.partD}.
In \path{IntegralHeckeForms.lean}, they are
\texttt{TwOK}, \texttt{GoodTw}, \texttt{IdxAt},
\texttt{TwistTheoremAtIndex}, and \texttt{admT4}.
In \path{TwGood.lean}, \texttt{twEps} fixes
$\ell^{-2}$ and \texttt{exists\_good\_of\_conditions}
uses precisely the existential selection above.
Those named definitions and the last theorem's
full proof were read directly; no Lean build is
claimed.  The finite-cover language uses only
the already audited SGA~1, Expos\'e~V, \S\S4 and 7.
